\textbf{Phần 1. Trắc nghiệm lựa chọn}
\setcounterpageref{ex}{0}

\begin{ex} %[5D7V41]
Một mảnh đất hình chữ nhật có diện tích $600\text{ m}^2$. Người ta dùng $35\%$ diện tích để xây nhà, $25\%$ diện tích làm sân chơi, phần còn lại trồng cây xanh. Diện tích đất trồng cây xanh là
\choice
{$180\text{ m}^2$}
{$210\text{ m}^2$}
{\True $240\text{ m}^2$}
{$360\text{ m}^2$}
\loigiai{
Tỉ số phần trăm diện tích trồng cây xanh là:
$100\% - (35\% + 25\%) = 40\%$.\\
Diện tích đất trồng cây xanh là:
$600 \times 40\% = 240\text{ (m}^2\text{)}$.\\
\textbf{Đáp án của câu hỏi là C}
}
\end{ex}

\begin{ex} %[5D7V41]
Một cuốn sách có giá bìa là $80\,000$ đồng. Nhân ngày hội sách, hiệu sách giảm giá $15\%$. Bạn Nam mua cuốn sách đó phải trả số tiền là
\choice
{$64\,000\text{ đồng}$}
{\True $68\,000\text{ đồng}$}
{$70\,000\text{ đồng}$}
{$72\,000\text{ đồng}$}
\loigiai{
Số tiền được giảm giá là:
$80\,000 \times 15\% = 12\,000\text{ (đồng)}$.\\
Số tiền bạn Nam phải trả là:
$80\,000 - 12\,000 = 68\,000\text{ (đồng)}$.\\
\textbf{Đáp án của câu hỏi là B}
}
\end{ex}

\begin{ex} %[5D7V41]
Một trang trại nuôi $1\,200$ con gà, trong đó số gà đẻ trứng chiếm $65\%$, còn lại là gà thịt. Số con gà thịt trang trại nuôi là
\choice
{$380\text{ con}$}
{\True $420\text{ con}$}
{$450\text{ con}$}
{$780\text{ con}$}
\loigiai{
Tỉ số phần trăm số gà thịt là:
$100\% - 65\% = 35\%$.\\
Số con gà thịt là:
$1\,200 \times 35\% = 420\text{ (con)}$.\\
\textbf{Đáp án của câu hỏi là B}
}
\end{ex}

\begin{ex} %[5D7V41]
Lương của một công nhân là $8\,500\,000$ đồng một tháng. Sau khi trừ đi $10{,}5\%$ các khoản đóng bảo hiểm, số tiền thực lĩnh người đó nhận được là
\choice
{$892\,500\text{ đồng}$}
{$7\,650\,000\text{ đồng}$}
{\True $7\,607\,500\text{ đồng}$}
{$7\,500\,000\text{ đồng}$}
\loigiai{
Số tiền đóng bảo hiểm là:
$8\,500\,000 \times 10{,}5\% = 892\,500\text{ (đồng)}$.\\
Số tiền thực lĩnh người đó nhận được là:
$8\,500\,000 - 892\,500 = 7\,607\,500\text{ (đồng)}$.\\
\textbf{Đáp án của câu hỏi là C}
}
\end{ex}

\begin{ex} %[5D7V41]
Một mảnh vườn hình chữ nhật có chiều dài $30\text{ m}$, chiều rộng bằng $70\%$ chiều dài. Diện tích của mảnh vườn đó là
\choice
{$420\text{ m}^2$}
{\True $630\text{ m}^2$}
{$840\text{ m}^2$}
{$900\text{ m}^2$}
\loigiai{
Chiều rộng mảnh vườn là:
$30 \times 70\% = 21\text{ (m)}$.\\
Diện tích mảnh vườn là:
$30 \times 21 = 630\text{ (m}^2\text{)}$.\\
\textbf{Đáp án của câu hỏi là B}
}
\end{ex}

\begin{ex} %[5D7V41]
Lượng nước trong hạt cà phê tươi chiếm $18\%$ khối lượng. Bác nông dân đem phơi $400\text{ kg}$ cà phê tươi. Lượng nước bốc hơi hoàn toàn khỏi số cà phê đó cân nặng bao nhiêu ki-lô-gam?
\choice
{$64\text{ kg}$}
{\True $72\text{ kg}$}
{$80\text{ kg}$}
{$328\text{ kg}$}
\loigiai{
Khối lượng nước bốc hơi là:
$400 \times 18\% = 72\text{ (kg)}$.\\
\textbf{Đáp án của câu hỏi là B}
}
\end{ex}

\begin{ex} %[5D7V41]
Một cửa hàng nhập về $250\text{ kg}$ cam. Ngày thứ nhất bán được $40\%$ số cam đó. Ngày thứ hai bán được $50\%$ số cam CÒN LẠI. Sau hai ngày cửa hàng còn lại bao nhiêu ki-lô-gam cam?
\choice
{$50\text{ kg}$}
{\True $75\text{ kg}$}
{$100\text{ kg}$}
{$125\text{ kg}$}
\loigiai{
Ngày thứ nhất bán được:
$250 \times 40\% = 100\text{ (kg)}$.\\
Số cam còn lại sau ngày thứ nhất là:
$250 - 100 = 150\text{ (kg)}$.\\
Ngày thứ hai bán được:
$150 \times 50\% = 75\text{ (kg)}$.\\
Số cam còn lại sau hai ngày là:
$150 - 75 = 75\text{ (kg)}$.\\
\textbf{Đáp án của câu hỏi là B}
}
\end{ex}

\begin{ex} %[5D7V41]
Một chiếc tivi có giá niêm yết $12\,000\,000$ đồng. Siêu thị giảm giá lần thứ nhất $10\%$. Sau đó siêu thị tiếp tục giảm thêm $5\%$ trên mức giá đã giảm. Giá bán cuối cùng của chiếc tivi là
\choice
{$10\,200\,000\text{ đồng}$}
{\True $10\,260\,000\text{ đồng}$}
{$10\,500\,000\text{ đồng}$}
{$10\,800\,000\text{ đồng}$}
\loigiai{
Giá sau lần giảm thứ nhất là:
$12\,000\,000 \times (100\% - 10\%) = 10\,800\,000\text{ (đồng)}$.\\
Giá sau lần giảm thứ hai là:
$10\,800\,000 \times (100\% - 5\%) = 10\,260\,000\text{ (đồng)}$.\\
\textbf{Đáp án của câu hỏi là B}
}
\end{ex}

\begin{ex} %[5D7V41]
Một lớp học có $40$ học sinh. Cuối học kì 1, có $30\%$ học sinh đạt mức Xuất sắc, $45\%$ học sinh đạt mức Giỏi, còn lại là mức Hoàn thành tốt. Số học sinh đạt mức Hoàn thành tốt của lớp là
\choice
{$8\text{ bạn}$}
{\True $10\text{ bạn}$}
{$12\text{ bạn}$}
{$15\text{ bạn}$}
\loigiai{
Tỉ số phần trăm học sinh Hoàn thành tốt là:
$100\% - (30\% + 45\%) = 25\%$.\\
Số học sinh Hoàn thành tốt là:
$40 \times 25\% = 10\text{ (bạn)}$.\\
\textbf{Đáp án của câu hỏi là B}
}
\end{ex}

\begin{ex} %[5D7V41]
Một bể chứa $2\,400\text{ l}$ nước. Buổi sáng người ta dùng hết $20\%$ lượng nước trong bể, buổi chiều dùng tiếp $25\%$ lượng nước ban đầu. Cả ngày người ta đã dùng hết bao nhiêu lít nước?
\choice
{$960\text{ l}$}
{$1\,000\text{ l}$}
{\True $1\,080\text{ l}$}
{$1\,200\text{ l}$}
\loigiai{
Cả ngày đã dùng số phần trăm lượng nước trong bể là:
$20\% + 25\% = 45\%$.\\
Số lít nước đã dùng cả ngày là:
$2\,400 \times 45\% = 1\,080\text{ (l)}$.\\
\textbf{Đáp án của câu hỏi là C}
}
\end{ex}

\begin{ex} %[5D7V41]
Bác An gửi tiết kiệm $60\,000\,000$ đồng với kì hạn $6$ tháng, lãi suất $0{,}5\%$ mỗi tháng. Sau $6$ tháng, số tiền lãi bác An nhận được là
\choice
{$300\,000\text{ đồng}$}
{$900\,000\text{ đồng}$}
{\True $1\,800\,000\text{ đồng}$}
{$3\,600\,000\text{ đồng}$}
\loigiai{
Tiền lãi trong một tháng là:
$60\,000\,000 \times 0{,}5\% = 300\,000\text{ (đồng)}$.\\
Sau 6 tháng số tiền lãi nhận được là:
$300\,000 \times 6 = 1\,800\,000\text{ (đồng)}$.\\
\textbf{Đáp án của câu hỏi là C}
}
\end{ex}

\begin{ex} %[5D7V41]
Một thửa ruộng thu hoạch được $3$ tấn thóc tươi. Sau khi phơi khô, khối lượng thóc hao hụt mất $8\%$. Số ki-lô-gam thóc khô thu được là
\choice
{$240\text{ kg}$}
{$2\,400\text{ kg}$}
{\True $2\,760\text{ kg}$}
{$2\,800\text{ kg}$}
\loigiai{
Đổi: $3\text{ tấn} = 3\,000\text{ kg}$.\\
Khối lượng thóc bị hao hụt là:
$3\,000 \times 8\% = 240\text{ (kg)}$.\\
Khối lượng thóc khô thu được là:
$3\,000 - 240 = 2\,760\text{ (kg)}$.\\
\textbf{Đáp án của câu hỏi là C}
}
\end{ex}

\textbf{Phần 2. Trắc nghiệm đúng sai}
\setcounterpageref{ex}{0}

\begin{ex} %[5D7V41]
Một siêu thị niêm yết giá một chiếc tủ lạnh là $8\,000\,000$ đồng. Nhân dịp ngày hội mua sắm, siêu thị giảm giá $15\%$.
\choiceTF
{\True Số tiền được giảm giá là $1\,200\,000$ đồng}
{\True Khách hàng cần trả $6\,800\,000$ đồng để mua chiếc tủ lạnh}
{\True Khách hàng thanh toán qua thẻ ngân hàng được hoàn lại thêm $5\%$ trên số tiền đã thanh toán, tương ứng được hoàn $340\,000$ đồng}
{Tổng số tiền khách hàng được giảm và hoàn lại vượt quá $1\,600\,000$ đồng}
\loigiai{
\begin{itemchoice}
\item Đúng. Số tiền giảm: $8\,000\,000 \times 15\% = 1\,200\,000\text{ (đồng)}$.
\item Đúng. Số tiền trả: $8\,000\,000 - 1\,200\,000 = 6\,800\,000\text{ (đồng)}$.
\item Đúng. Số tiền hoàn: $6\,800\,000 \times 5\% = 340\,000\text{ (đồng)}$.
\item Sai. Tổng tiền giảm và hoàn: $1\,200\,000 + 340\,000 = 1\,540\,000\text{ (đồng)} < 1\,600\,000\text{ đồng}$.
\end{itemchoice}
}
\end{ex}

\begin{ex} %[5D7V41]
Một mảnh đất hình chữ nhật có chu vi $120\text{ m}$, chiều rộng bằng $60\%$ chiều dài.
\choiceTF
{\True Nửa chu vi mảnh đất là $60\text{ m}$}
{\True Chiều dài của mảnh đất là $37{,}5\text{ m}$}
{\True Chiều rộng của mảnh đất là $22{,}5\text{ m}$}
{Diện tích của mảnh đất là $800\text{ m}^2$}
\loigiai{
\begin{itemchoice}
\item Đúng. Nửa chu vi: $120 : 2 = 60\text{ (m)}$.
\item Đúng. Chiều dài: $60 : (1 + 0{,}6) = 37{,}5\text{ (m)}$.
\item Đúng. Chiều rộng: $37{,}5 \times 0{,}6 = 22{,}5\text{ (m)}$.
\item Sai. Diện tích: $37{,}5 \times 22{,}5 = 843{,}75\text{ (m}^2\text{)} 
eq 800\text{ m}^2$.
\end{itemchoice}
}
\end{ex}

\textbf{Phần 3. Trắc nghiệm trả lời ngắn}
\setcounterpageref{ex}{0}

\begin{ex} %[5D7V41]
Tìm giá trị của $25\%$ của số $360$.
\shortans{90}
\loigiai{
Giá trị cần tìm là:
$360 \times 25\% = 90$.
}
\end{ex}

\begin{ex} %[5D7V41]
Một lớp học có $40$ học sinh, trong đó có $45\%$ là học sinh nam. Hỏi lớp học đó có bao nhiêu học sinh nữ?
\shortans{22}
\loigiai{
Tỉ số phần trăm học sinh nữ là:
$100\% - 45\% = 55\%$.\\
Số học sinh nữ của lớp là:
$40 \times 55\% = 22\text{ (học sinh)}$.
}
\end{ex}

\begin{ex} %[5D7V41]
Một mảnh vườn có diện tích $500\text{ m}^2$. Người ta dùng $36\%$ diện tích để trồng cây ăn quả. Diện tích đất trồng cây ăn quả là bao nhiêu mét vuông?
\shortans{180}
\loigiai{
Diện tích trồng cây ăn quả là:
$500 \times 36\% = 180\text{ (m}^2\text{)}$.
}
\end{ex}

\begin{ex} %[5D7V41]
Một thùng chứa $60\text{ l}$ dầu. Người ta đã lấy đi $40\%$ số dầu trong thùng. Hỏi trong thùng còn lại bao nhiêu lít dầu?
\shortans{36}
\loigiai{
Số dầu còn lại chiếm:
$100\% - 40\% = 60\%$.\\
Số lít dầu còn lại là:
$60 \times 60\% = 36\text{ (l)}$.
}
\end{ex}

\begin{ex} %[5D7V41]
Một cửa hàng có $150$ bộ quần áo. Buổi sáng bán được $20\%$ số quần áo. Buổi chiều bán được $30\%$ số quần áo ban đầu. Cả ngày cửa hàng đã bán được bao nhiêu bộ quần áo?
\shortans{75}
\loigiai{
Cả ngày bán được số phần trăm là:
$20\% + 30\% = 50\%$.\\
Số bộ quần áo bán được là:
$150 \times 50\% = 75\text{ (bộ)}$.
}
\end{ex}

\begin{ex} %[5D7V41]
Giá niêm yết của một chiếc xe đạp là $2\,000\,000$ đồng. Nhân dịp khai trương, cửa hàng giảm giá $15\%$. Giá bán của chiếc xe đạp sau khi giảm giá là bao nhiêu nghìn đồng?
\shortans{1700}
\loigiai{
Giá bán sau khi giảm là:
$2\,000\,000 \times (100\% - 15\%) = 1\,700\,000\text{ (đồng)}$.\\
Đổi: $1\,700\,000\text{ đồng} = 1\,700\text{ nghìn đồng}$.
}
\end{ex}
